\input{../common}
\everymath{\displaystyle}
\begin{document}
  %<*content>
  \development{algebra}{situationscomplexes}{Situations complexes}

  \summary{Compilés par M.Mbodji}



\begin{exercice}
	Sophie est une élève en classe de 6\up{ième} a eu le paludisme à la fin du mois de Septembre pendant les vacances et souhaite rentrer à la maison après 30 heures d'hospitalisation pour préparer sa rentrée scolaire. Un médecin lui a injecté un médicament par voie intraveineuse à l'hôpital, et dans les heures qui suivent la substance est éliminée par les reins. La quantité (en mg)  présente dans le sang à l'instant  (en heures) a été mesurée par des prises de sang toutes les six heures comme indique le tableau ci-dessous. Sophie peut quitter l'hôpital si la substance est totalement éliminée de son organisme.
\begin{center}\vspace*{-0.5em}
\renewcommand{\arraystretch}{1.5}
\begin{tabular}{|c|c|c|c|c|c|}
\hline
   $ x_i $   & 0 & 6 &  12&  18&  24 \\
\hline
 $ y_i $   & 29,7 & 22,5 & 16,5 & 11,7 & 9  \\
\hline
\end{tabular}
\end{center}


Sophie  peut-elle sereinement préparer sa rentrée scolaire ?  
\end{exercice}

\begin{exercice}
Le plan complexe est rapporté à un repère  orthonormé $ \ouv $. L'unité graphique est le mètre.
\medskip

$ \bullet $  \textbf{M.DIOP} a un terrain de forme rectangulaire  dont les dimensions $x$ et $y$ sont tels que:
\medskip

\hspace*{3cm} $ (2+3\i)z+(1-3\i)\overline{z} =6+3\i$  \qquad  $ z=x+\i y $ 

\medskip

\noindent Il voudrait construire sur ce terrain une école, et pour cela il a besoin de recouvrir  toute la superficie de ce terrain avec des carreaux. Le carton de carreaux coûte 14 000 FCFA et peut recouvrir une superficie de 5$m^2$.


\medskip


$ \bullet $  \textbf{M.NGOM} quant à lui  a un terrain de forme triangulaire  dont un sommet est repéré par son affixe  $ 2+3\i $ et les deux autres  sommets  ont des affixes   solutions de l'équation :

\medskip
\hspace*{3cm}  $ z^2 + (2+3\i)z -2(1-2\i)=0$

\medskip
\noindent Il souhaite le clôturer à l'aide d'un grillage dont le mètre coûte 1500 FCFA.

\medskip

$ \bullet $ Le terrain que  \textbf{M.NDIAYE}  possède est situé en plein  quartier administratif  dont la forme est celle des points $ M $ d'affixes $ z\neq -1+2\i y $ tel que \; $ \dfrac{z-7+4\i}{z+1-2\i}$ soit un imaginaire pur. Il souhaite l'hypothéquer  avec une voiture dont la valeur est estimée à 1 170 000 FCFA. Sachant que son terrain a une valeur de $ 15000$  F CFA le mètre carré.

\medskip

\textbf{Votre travail en tant qu'élève de TS2, consiste à résoudre les tâches suivantes en justifiant votre démarche par un raisonnement bien détaillé}.

\medskip

\textbf{Tâches:}
\begin{enumerate}
\item Quelle somme doit dépenser  \textbf{M.NGOM} pour clôturer son jardin ?
\item Déterminer une estimation du montant nécessaire pour l'achat des carreaux devant recouvrir entièrement le terrain de  \textbf{M.DIOP}.
\item \textbf{M.NDIAYE}  réussira- t-il à être propriétaire de cette voiture? 
\end{enumerate}
\end{exercice}
\begin{exercice}
\noindent
Le plan  est rapporté à un repère  orthonormé $ \ouv $. L'unité graphique  est le  hectomètre.
\medskip

\noindent\textbf{$  \bullet$}\; 
 On appelle couverture réseau de rayon $ r $, l'aire de disque de rayon $ r $ et de centre le point de fixation d'un pylône.

\noindent La couverture réseau d'une
société de téléphonie dans une ville  est donnée par l'ensemble  $\mathcal{C}$ des points $M$ d'affixe $z$ tel que:
 $ \bigl|z -2-3\i\bigr| \leq 3$. 
 
 \smallskip
 \noindent Par ailleurs, dans cette ville, un quartier est délimité par un espace de forme triangulaire
dont les sommets sont les images des  solutions de l'équation (E) :\; 
  $ z^3 -(5+7\i)z^2 - (4-25\i)z -12\i+30=0$, l'un des sommets a pour coordonnées \, $ (0  \; ; \;  2) $. 


\bigskip

\noindent \textbf{$  \bullet$}\;  Monsieur NGOM dispose d'un terrain ayant la forme d'un triangle dont les sommets sont 
les  images des solutions de  l'équation (E'):\; $ z^3=8\i $,\; l'un des sommets est repéré par son affixe $ -2\i $. 


\noindent  Il souhaite clôturer son terrain à l'aide de cinq tours de
grillage dont le mètre coûte 2 200 fcfa.




\medskip

\noindent  \textbf{\textit{Votre travail en tant qu'élève de TS2, consiste à résoudre les tâches suivantes en justifiant votre démarche par un raisonnement bien détaillé}}.

\medskip

\textbf{\underline{Tâches:}}
\begin{enumerate}
\item Le quartier est-il entièrement couvert par le réseau de cette société ?
\item Quelle somme doit dépenser  M. NGOM  pour clôturer son terrain?
\end{enumerate}
\end{exercice}

\begin{exercice}
 On considère dans le plan complexe muni du repère orthonormal $ \ouv $, les points $A $ , $ B$ , $ C $  et $ D $ d'affixes respectives $ z_A=3+\i $,\;  $z_B=1+3\i $ \; , \;  $ z_C=1- \i $ \; et \; $ z_D=\overline{z_A} $.

\medskip
 
  Le jardin potager de M. Mbaye  est formé du quadrilatère $ ABCD $ qu'il voudrait clôturer  par un fil barbelé en laissant une porte de $ 0.8 $ mètre. Le rouleau de 5m de
ce fil lui est vendu à 3500 FCFA. (\textit{On prendra 1 m pour unité})


 Combien va-t-il dépenser pour clôturer son jardin. 


\end{exercice}

\begin{exercice}
Une usine fabrique et commercialise des sachets de
poudre de jus fruits. Sa capacité journalière de
production est comprise entre 1 000 et 3 000
sachets. On suppose que toute la production est
commercialisée. Une étude a révélé que le bénéfice
journalier, exprimé en millions de francs CFA,
réalisé pour la production et la vente de milliers
de sachets est modélisé sur l'intervalle
[1 ; 3] par la fonction B définie par :\;
$ B(x)= -\frac{1}{2}x^2+x+2+2\ln x$\\
Le Directeur de l'usine veut accroître le Bénéfice
de l'entreprise. N'ayant pas de personnel qualifié,
il te demande le nombre de sachets à produire en
un jour, à l'unité près, pour que l'entreprise réalise
un bénéfice maximal.
\end{exercice}
\begin{exercice}
 Un enseignant de Mathématiques pour récompenser le meilleur élève de sa classe après une évaluation convient avec ses élèves d'un jeu qui consiste à tirer  au hasard et simultanément trois  boules dans une urne   contenant 7 boules blanches et 5 boules rouges toutes indiscernables au toucher.. 
 

 Chaque   boule blanche tirée  rapporte 1 000 F et
chaque boule rouge rapporte 2 000 F.


\medskip
\noindent Avant de tirer, l'élève Pierre, le meilleur,  
exprime son
désir d'obtenir   5 000F  ou plus.


\medskip
\noindent Pour lui, il aura dans ce cas, au moins une chance sur deux
d'avoir la somme  souhaitée.

\medskip
\noindent  Pierre a t-il  raison dans son affirmation ?  Justifier.
\end{exercice}
\begin{exercice}
Une étude sur les rendements de cultures pluviales de riz dans un pays de l'Afrique de l'Ouest révèle que le taux de diminution de la production est proportionnel à la production en chaque année.

\medskip

Initialement la production est de 500 000 tonnes et elle diminue de 400 000 tonnes en cinq ans.

\medskip

Selon la conclusion de cette étude ; d'ici à plus de dix ans,\, 50\% de la production initiale diminuerait si les conditions climatiques restent inchangées.

\medskip
Donne ton avis sur cette conclusion de cette étude.
\end{exercice}
\begin{exercice}
Pour bitumer une route de banlieue longue de 800 mètres, un maire lance un appel d’offre pour
désigner le moins disant. Deux entreprises A et B ont soumissionné et estiment le coût de bitumage
comme suit.
\medskip

\noindent Pour l’entreprise A, le premier décamètre coûte 100.000 F, et chaque décamètre suivant coûtera
1.000 F moins cher que le précédent.
\medskip

\noindent Pour l’entreprise B, le premier hectomètre coûte 500.000 F, et chaque hectomètre suivant coûtera
5\% plus cher que le précédent.

\medskip

\noindent \textbf{Tache 1 }: on note $A_n$ le prix du $n-ième$ décamètre proposé par l’entreprise A et $B_n$ le prix du $n-ième$ hectomètre proposé par l’entreprise B.

\noindent Exprimer $A_n$ et $B _n$ en fonction de $n$ et déterminer la nature de
chacune des suites $(A_n )$ et $(B_n )$.


\medskip
\noindent\textbf{Tache 2} : Calculer le coût de bitumage proposé par chaque entreprise et en déduire la moins disant
(celle qui propose un moindre coût)


\end{exercice}
\begin{exercice}



Le tableau ci-dessous donne le nombre total d’adhérents au club littéraire d’un lycée au cours de l’année civile 2020.

\begin{center}
\begin{tabular}{|c|c|c|c|c|c|c|c|c|c|c|c|c|}
\hline
Mois & Jan & Fév & Mar & Avr & Mai & Juin & Juil & Août & Sept & Oct & Nov & Déc \\
\hline
Rang $x_i$ & 1 & 2 & 3 & 4 & 5 & 6 & 7 & 8 & 9 & 10 & 11 & 12 \\\hline
Nombre $y_i$ & 1100 & 1160 & 1220 & 1370 & 1620 & 1550 & 1600 & 1500 & 1790 & 1940 & 2060 & 1980 \\
\hline
\end{tabular}
\end{center}

Une Organisation Non Gouvernementale promet d’octroyer une aide financière considérable au club si le nombre d’adhérents dépasse les 3000 élèves. L’élève de la Terminale S2 qui dirige le club désire connaître la date à laquelle ce don pourra se faire. Il te sollicite pour l’aider. Détermine la date (mois et année) probable de la réception de ce don.

\end{exercice}
\begin{exercice}
La grippe est une infection respiratoire contagieuse due aux virus \emph{influenza}. Elle se propage facilement lorsqu’une personne tousse ou éternue. Les médecins décident de lutter contre la propagation de la maladie dans la population en mettant au point un test pour le dépistage de cette grippe. Ils considèrent que le test est fiable lorsqu’au moins 99 personnes sur 100 ayant un test positif sont réellement malades. 

Le laboratoire faisant ce test fournit les caractéristiques suivantes :

\begin{itemize}
    \item Une personne malade a 49 chances sur 50 d’être testée positive.
    \item Une personne non malade a une chance sur 1000 d’être testée positive.
\end{itemize}

Les médecins procèdent à un test de dépistage systématique dans une population « cible ». On note $p$, la proportion de personnes atteintes dans la population « cible ».

Soit $f$ la proportion d’individus susceptibles d’être malades sachant qu’ils sont déclarés positifs par le test.

\begin{enumerate}
    \item Montrer que $f$ s’exprime en fonction de $p$ par : 
  $  
    f(p) = \dfrac{980p}{979p + 1}
   $
    \item À partir de quelle proportion de personnes atteintes dans la population jugent-ils que le test est fiable ?
\end{enumerate}
\end{exercice}


\begin{exercice}
En début d’année scolaire 2025, un élève de terminale qui habite à trois kilomètres de son établissement décide d’acheter un scooter et un réveil pour ne pas arriver en retard en classe.\\
On sait qu’il peut être victime de deux situations indépendantes, à chaque matin de classe :

\begin{itemize}
  \item $R$ : il n’entend pas son réveil sonner ;
  \item $S$ : son scooter, mal entretenu, tombe en panne.
\end{itemize}

\begin{description}
  \item Il a observé que chaque jour de classe, il y a 2 possibilités sur 10 que la situation $R$ se produise et 5 possibilités sur 100 que la situation $S$ se produise.

Lorsqu’au moins une de ces situations se produit, l’élève arrive en retard au lycée, sinon il est à l’heure.

\item Au cours d'une  semaine, cet élève se rend 5 fois de suite en classe. On admet que le fait qu’il soit en retard un jour de classe donné n’influe pas sur le fait qu’il le soit ou non les jours suivants.
\item 
Quelle est la chance qu’il soit à l’heure au moins quatre jours sur les cinq durant une semaine ?\\
(Donner les résultats sous forme de fraction irréductible.)
\end{description}
\end{exercice}

\begin{exercice}
Une association prévoit d'organiser une cérémonie dans la place publique
le 5 Août 2024 mais il y a la menace de l'hivernage vu que le mois d'Août
est très pluvieux.

Pour s'assurer de la tenue de l'évènement, le bureau va demander les
services d'un institut météorologique, ce dernier donne les informations
suivantes :
\begin{itemize}
\item La probabilité qu'il pleuve le premier Août 2024 est de  $ \frac{1}{4} $
\item S'il a plu un jour dans le mois la probabilité qu'il pleuve le jour suivant est $ \frac{1}{2} $
\item S’il n'a pas plu un jour dans le mois la probabilité qu’il ne pleuve pas
le jour suivant est $ \frac{1}{5} $.
\end{itemize}
Le bureau a décidé d'entamer les préparatifs de la cérémonie que si la
probabilité qu'il pleuve le jour de la cérémonie est inférieure à 0,5 sinon il va
changer la date.

\noindent On note par $ A_n $ l'évènement << Il a plu le n-ieme du mois d'Août >>.

\noindent Et  $ p_n $ la probabilité de l'évènement \; $ A_n $\, soit   $p_n=P(A_n)$.
\begin{enumerate}
%\item Calculer $ p_2 $ .
%\item  Donner \; $ P \paren{A_{n+1} /A_{n}} $\;  et \; $ P \paren{A_{n+1} /\overline{A_{n}}} $.
%\item Montrer que \; $ p_{n+1}=-\frac{3}{10}p_n +\frac{4}{5}$
\item En utilisant la suite géométrique  $(U_  n)$ définie par  $U_n=p_n -\frac{8}{13}$, exprimer 
 $p_n$ en fonction de $n$.

\item L’association va-t-elle entamer les préparatifs pour l’organisation de la
cérémonie ? Justifier la réponse.
\end{enumerate} 
\end{exercice}
\begin{exercice}  
Une usine fabrique et commercialise des sachets de
poudre de jus fruits. Sa capacité journalière de
production est comprise entre 1 000 et 3 000
sachets. On suppose que toute la production est
commercialisée. Une étude a révélé que le bénéfice
journalier, exprimé en millions de francs CFA,
réalisé pour la production et la vente de milliers
de sachets est modélisé sur l'intervalle
[1 ; 3] par la fonction B définie par :\\
$ B(x)= -\frac{1}{2}x^2+X+2+2\ln x$\\
Le Directeur de l'usine veut accroître le Bénéfice
de l'entreprise. N'ayant pas de personnel qualifié,
il te demande le nombre de sachets à produire en
un jour, à l'unité près, pour que l'entreprise réalise
un bénéfice maximal.
\end{exercice} 

\begin{exercice}


Ami et Rémi   ont commencé à travailler pour différentes entreprises le  1\up{er}  janvier 2011.

\begin{enumerate}
\item[$  \bullet$]  Le salaire annuel de départ de  Rémi  était de $ 450000 $ fcfa
 et son salaire annuel augmente de 2\% le   1\up{er}  janvier de chaque année après  2011.
\medskip


\item[$  \bullet$] Le salaire  annuel de  Ami  est basé sur une évaluation annuelle des performances. Son salaire  annuel pour les années 2011, 2013, 2014, 2018, 2022 est indiqué dans le tableau suivant.
 \begin{center}
 \renewcommand{\arraystretch}{1.2}
 \setlength{\arrayrulewidth}{1.2pt}
 %\arrayrulecolor{blue}
 \begin{tabular}{|c|c|c|c|c|c|}
 \hline
 Année (x) &2011&2013&2014&2018&2022\\ \hline
 Salaire annuel(y) (en 10000 Fcfa) & 45& 47,2&48,5&53&57\\\hline
 
 \end{tabular} 
 \end{center}
\end{enumerate}

\textbf{ Votre travail en tant qu'élève de TS2, consiste à résoudre la tâche suivante en justifiant votre démarche par un raisonnement bien détaillé.}

\medskip

\textbf{Tâche:}
 En supposant que le salaire annuel de Ami suit cette évolution, comparez les salaires 
de Ami et  de Rémi au cours de l’année 2021.
   
\end{exercice}

\end{document}